\documentclass[11pt]{article}

\usepackage{arxiv}
\usepackage[spanish]{babel}
\usepackage[utf8]{inputenc}
\usepackage[T1]{fontenc}
\usepackage{amsmath}
\usepackage{amsfonts}
\usepackage{booktabs}
\usepackage{graphicx}
\usepackage{hyperref}
\usepackage{microtype}
\usepackage{nicefrac}
\usepackage{url}

\graphicspath{{./figures/mm1/}{./figures/inventory/}}
\renewcommand{\abstractname}{Resumen}
\renewcommand{\keywordname}{\bfseries \emph{Palabras clave}}
\renewcommand{\headeright}{Preprint}
\renewcommand{\undertitle}{Preprint}
\renewcommand{\shorttitle}{TP 3 - M/M/1 e inventario}
\hypersetup{
    colorlinks=true,
    linkcolor=blue,
    urlcolor=blue,
    citecolor=blue
}

\title{TP 3 - Simulacion de un modelo M/M/1 e inventario}
\author{Simulacion computacional}
\date{\today}

\begin{document}

\maketitle
\begin{abstract}
Este informe presenta el desarrollo de dos modelos del trabajo practico: un sistema de colas M/M/1 y un modelo de inventario con politica de revision continua $(Q,r)$. Para M/M/1 se documentan los supuestos teoricos, el diseno experimental, la simulacion por eventos discretos en Python y la comparacion contra resultados estacionarios cuando corresponde. Los casos de cola infinita con $\rho \geq 1$ se conservan como simulaciones de horizonte finito, pero no se comparan contra teoria estacionaria porque el modelo no posee regimen estable en esas condiciones. Para inventario se comparan tres politicas bajo demanda diaria Poisson, ventas perdidas y costos de orden, mantenimiento y faltante, incorporando una referencia esperada simplificada como punto de contraste. La comparacion con AnyLogic se integrara cuando esten disponibles los resultados del modelo desarrollado en paralelo.
\end{abstract}
\keywords{simulacion \and eventos discretos \and M/M/1 \and colas finitas \and inventario}

\section{Introduccion}

El trabajo practico estudia sistemas mediante simulacion computacional y comparacion con resultados teoricos o esperados cuando estos existen. Los modelos analizados son una cola M/M/1, implementada mediante una simulacion de eventos discretos en Python, y un sistema de inventario con politica de revision continua $(Q,r)$, simulado por periodos diarios.

Para el modelo de colas se busca observar como cambia la congestion al variar la relacion entre tasa de arribo y tasa de servicio. Ademas, se analizan colas con capacidad finita para medir la probabilidad de denegacion de servicio cuando el sistema no puede aceptar nuevos clientes.

Los parametros principales se leen desde \texttt{config/default\_parameters.json}, lo que permite modificar la configuracion experimental sin cambiar el codigo fuente.

La implementacion en AnyLogic se esta desarrollando en paralelo y se incorporara manualmente al informe final cuando existan resultados exportados comparables.

\section{Marco teorico del modelo M/M/1}

\subsection{Supuestos del modelo}

El modelo M/M/1 representa un sistema de espera con una sola estacion de servicio. La primera M indica que los tiempos entre arribos siguen una distribucion exponencial, equivalente a un proceso de Poisson para la cantidad de arribos. La segunda M indica que los tiempos de servicio tambien son exponenciales. El numero 1 indica que existe un unico servidor.

En este trabajo se asume disciplina FIFO, es decir, los clientes son atendidos en el mismo orden en el que llegan. Las tasas se expresan en clientes por hora:

\begin{itemize}
    \item $\lambda$: tasa media de arribos.
    \item $\mu$: tasa media de servicio.
    \item $\rho = \lambda / \mu$: factor de carga o utilizacion teorica del sistema con cola infinita.
\end{itemize}

\subsection{Cola infinita}

Para una cola infinita, el sistema solo tiene regimen estacionario si:

\begin{equation}
    \rho < 1.
\end{equation}

Cuando esta condicion se cumple, las medidas teoricas usadas para validar la simulacion son:

\begin{align}
    L &= \frac{\rho}{1-\rho}, \\
    L_q &= \frac{\rho^2}{1-\rho}, \\
    W &= \frac{1}{\mu-\lambda}, \\
    W_q &= \frac{\lambda}{\mu(\mu-\lambda)}, \\
    U &= \rho.
\end{align}

Donde $L$ es el numero promedio de clientes en el sistema, $L_q$ es el numero promedio de clientes en cola, $W$ es el tiempo promedio en el sistema, $W_q$ es el tiempo promedio de espera en cola y $U$ es la utilizacion del servidor.

Para la distribucion de longitud de cola se usa la relacion entre el numero de clientes en sistema $N$ y el numero de clientes esperando $Q$. Si $Q=0$, el sistema puede estar vacio o tener un cliente en servicio. Por lo tanto:

\begin{align}
    P(Q=0) &= P(N=0) + P(N=1), \\
    P(Q=q) &= P(N=q+1), \quad q \geq 1.
\end{align}

Como en M/M/1 estable $P(N=n)=(1-\rho)\rho^n$, queda:

\begin{align}
    P(Q=0) &= (1-\rho)(1+\rho), \\
    P(Q=q) &= (1-\rho)\rho^{q+1}, \quad q \geq 1.
\end{align}

Si $\rho \geq 1$, la cola infinita no tiene valores estacionarios finitos. En esos casos la simulacion por horizonte finito puede generar metricas observadas, pero no corresponde compararlas contra las formulas estacionarias anteriores.

\subsection{Cola finita}

Para estudiar denegacion de servicio se usa una cola con capacidad finita. En el codigo y en las tablas, $B$ representa la cantidad maxima de clientes que pueden esperar en cola. Como ademas puede haber un cliente en servicio, la capacidad total del sistema es:

\begin{equation}
    K = B + 1.
\end{equation}

El estado $n$ representa la cantidad de clientes dentro del sistema, con $0 \leq n \leq K$. Las probabilidades estacionarias para M/M/1/K son:

\begin{equation}
    P_n =
    \begin{cases}
        \frac{(1-\rho)\rho^n}{1-\rho^{K+1}}, & \rho \neq 1, \\
        \frac{1}{K+1}, & \rho = 1.
    \end{cases}
\end{equation}

La probabilidad de denegacion es la probabilidad de que el sistema este lleno al momento de un arribo:

\begin{equation}
    P_{den} = P_K.
\end{equation}

La tasa efectiva de clientes aceptados queda:

\begin{equation}
    \lambda_{\mathrm{eff}} = \lambda(1-P_{den}).
\end{equation}

Con estas probabilidades, las medidas teoricas se calculan como:

\begin{align}
    L &= \sum_{n=0}^{K} n P_n, \\
    L_q &= \sum_{n=0}^{K} \max(n-1,0) P_n, \\
    W &= \frac{L}{\lambda_{\mathrm{eff}}}, \\
    W_q &= \frac{L_q}{\lambda_{\mathrm{eff}}}, \\
    U &= 1-P_0.
\end{align}

A diferencia de la cola infinita, el modelo con capacidad finita tiene regimen estacionario tambien para $\rho \geq 1$, porque los arribos excedentes se rechazan cuando el sistema esta lleno.

\section{Metodologia de simulacion M/M/1}

\subsection{Diseno experimental}

La unidad de tiempo elegida para el modelo de colas es la hora. Se fija una tasa de servicio base:

\begin{equation}
    \mu = 10 \text{ clientes/hora}.
\end{equation}

Las tasas de arribo se definen como fracciones de $\mu$:

\begin{equation}
    \lambda \in \{0.25\mu, 0.50\mu, 0.75\mu, 1.00\mu, 1.25\mu\}.
\end{equation}

Por lo tanto, los valores de $\rho$ evaluados son $0.25$, $0.50$, $0.75$, $1.00$ y $1.25$. Para cada valor se simula una cola infinita y colas finitas con:

\begin{equation}
    B \in \{0,2,5,10,50\}.
\end{equation}

El horizonte de cada corrida es de $10000$ horas y se ejecutan $10$ replicas por experimento. Las semillas se derivan de una semilla base registrada, para que los resultados sean reproducibles.

\subsection{Logica de eventos discretos}

La simulacion mantiene un reloj y procesa dos tipos de eventos:

\begin{itemize}
    \item llegada de un cliente;
    \item finalizacion de un servicio.
\end{itemize}

En cada llegada se agenda el siguiente arribo. Si el servidor esta libre, el cliente inicia servicio inmediatamente. Si el servidor esta ocupado y existe lugar en la cola, el cliente espera. Si la cola esta llena, el cliente se rechaza y se incrementa el contador de denegaciones.

Los tiempos entre arribos y los tiempos de servicio se generan con variables exponenciales:

\begin{align}
    T_a &\sim \mathrm{Exp}(\lambda), \\
    T_s &\sim \mathrm{Exp}(\mu).
\end{align}

\subsection{Metricas registradas}

Las medidas temporales se calculan integrando el estado del sistema a lo largo del horizonte simulado. Si $T$ es el tiempo total de simulacion:

\begin{align}
    \widehat{L} &= \frac{1}{T}\int_0^T N(t)\,dt, \\
    \widehat{L_q} &= \frac{1}{T}\int_0^T Q(t)\,dt, \\
    \widehat{U} &= \frac{\text{tiempo ocupado del servidor}}{T}.
\end{align}

Para cada cliente atendido tambien se registra el tiempo de espera en cola y el tiempo total dentro del sistema. A partir de esos datos se estiman:

\begin{align}
    \widehat{W} &= \text{tiempo promedio en sistema}, \\
    \widehat{W_q} &= \text{tiempo promedio en cola}.
\end{align}

La probabilidad de denegacion se estima como:

\begin{equation}
    \widehat{P}_{den} =
    \frac{\text{clientes rechazados}}{\text{clientes generados}}.
\end{equation}

\subsection{Resumen estadistico}

Para cada combinacion de parametros se calcula la media de las $10$ replicas, el desvio estandar muestral y un intervalo de confianza del $95\%$:

\begin{equation}
    \bar{X} \pm t_{0.975,9}\frac{s}{\sqrt{10}}.
\end{equation}

Cuando existe valor teorico estacionario, se reporta ademas el error porcentual:

\begin{equation}
    \text{Error \%} =
    100 \frac{\bar{X} - X_{\mathrm{teorico}}}{X_{\mathrm{teorico}}}.
\end{equation}

Los experimentos de cola infinita con $\rho \geq 1$ se conservan porque forman parte de las tasas pedidas, pero se etiquetan como casos sin regimen estacionario. En consecuencia, sus metricas se interpretan como resultados de horizonte finito y no como aproximaciones a una teoria estacionaria.

\section{Resultados M/M/1}

Los resultados de Python se exportaron a \texttt{results/mm1/experiments/}. El archivo \texttt{mm1\_runs.csv} contiene una fila por corrida, mientras que \texttt{mm1\_summary.csv} y \texttt{mm1\_summary\_long.csv} contienen resumenes por experimento.

\subsection{Cola infinita estable}

Para $\rho=0.25$, $\rho=0.50$ y $\rho=0.75$, la cola infinita tiene regimen estacionario y puede compararse contra las formulas teoricas. La Tabla~\ref{tab:mm1-infinite-steady} muestra que las medias simuladas quedan cerca de los valores teoricos.

% Generado desde results/mm1/experiments/mm1_summary.csv
\begin{table}[htbp]
\centering
\caption{Comparacion teorica y simulada para cola infinita estable.}
\label{tab:mm1-infinite-steady}
\begin{tabular}{llrrrr}
\toprule
$\rho$ & Metrica & Sim. & IC 95\% & Teoria & Error \% \\
\midrule
0.25 & $L$ & 0.3332 & 0.0026 & 0.3333 & -0.04 \\
 & $L_q$ & 0.0829 & 0.0015 & 0.0833 & -0.53 \\
 & $W$ & 0.1329 & 0.0006 & 0.1333 & -0.31 \\
 & $W_q$ & 0.0331 & 0.0005 & 0.0333 & -0.81 \\
 & $\rho_{serv}$ & 0.2503 & 0.0013 & 0.2500 & 0.13 \\
\addlinespace
0.50 & $L$ & 1.0068 & 0.0074 & 1.0000 & 0.68 \\
 & $L_q$ & 0.5052 & 0.0060 & 0.5000 & 1.05 \\
 & $W$ & 0.2009 & 0.0013 & 0.2000 & 0.43 \\
 & $W_q$ & 0.1008 & 0.0011 & 0.1000 & 0.80 \\
 & $\rho_{serv}$ & 0.5015 & 0.0019 & 0.5000 & 0.31 \\
\addlinespace
0.75 & $L$ & 3.0205 & 0.0738 & 3.0000 & 0.68 \\
 & $L_q$ & 2.2705 & 0.0706 & 2.2500 & 0.91 \\
 & $W$ & 0.4025 & 0.0089 & 0.4000 & 0.63 \\
 & $W_q$ & 0.3026 & 0.0087 & 0.3000 & 0.85 \\
 & $\rho_{serv}$ & 0.7500 & 0.0039 & 0.7500 & 0.00 \\
\bottomrule
\end{tabular}
\end{table}


La Tabla~\ref{tab:mm1-queue-probabilities} reporta la probabilidad de observar $q$ clientes esperando en cola. Esta salida proviene del tiempo acumulado por longitud de cola durante la simulacion y se compara contra la distribucion estacionaria correspondiente.

% Generado desde results/mm1/experiments/mm1_queue_probabilities.csv
\begin{table}[htbp]
\centering
\caption{Probabilidad de longitud de cola para cola infinita estable.}
\label{tab:mm1-queue-probabilities}
\begin{tabular}{lccc}
\toprule
$q$ & $\rho=0.25$ & $\rho=0.50$ & $\rho=0.75$ \\
\midrule
$0$ & \shortstack{0.9380\\(0.9375)} & \shortstack{0.7489\\(0.7500)} & \shortstack{0.4372\\(0.4375)} \\
$1$ & \shortstack{0.0466\\(0.0469)} & \shortstack{0.1247\\(0.1250)} & \shortstack{0.1402\\(0.1406)} \\
$2$ & \shortstack{0.0115\\(0.0117)} & \shortstack{0.0628\\(0.0625)} & \shortstack{0.1049\\(0.1055)} \\
$3$ & \shortstack{0.0029\\(0.0029)} & \shortstack{0.0316\\(0.0312)} & \shortstack{0.0786\\(0.0791)} \\
$4$ & \shortstack{0.0008\\(0.0007)} & \shortstack{0.0161\\(0.0156)} & \shortstack{0.0588\\(0.0593)} \\
$5$ & \shortstack{0.0002\\(0.0002)} & \shortstack{0.0079\\(0.0078)} & \shortstack{0.0444\\(0.0445)} \\
\bottomrule
\end{tabular}
\par\smallskip
\begin{minipage}{0.95\linewidth}
\centering\footnotesize
Cada celda muestra simulacion media y, entre parentesis, valor teorico de $P(Q=q)$.
\end{minipage}
\end{table}


\subsection{Cola infinita sin regimen estacionario}

Para $\rho=1.00$ y $\rho=1.25$, la cola infinita no tiene estado estacionario. Las corridas se mantienen para cumplir el barrido experimental, pero sus valores no deben interpretarse como una convergencia teorica.

% Generado desde results/mm1/experiments/mm1_summary.csv
\begin{table}[htbp]
\centering
\caption{Corridas con cola infinita sin regimen estacionario.}
\label{tab:mm1-infinite-unstable}
\begin{tabular}{rrrrrl}
\toprule
$\rho$ & $L$ sim. & $L_q$ sim. & $W$ sim. & $W_q$ sim. & Referencia \\
\midrule
1.00 & 186.3341 & 185.3393 & 18.6264 & 18.5265 & No estacionaria \\
1.25 & 12541.8921 & 12540.8922 & 1003.8583 & 1003.7582 & No estacionaria \\
\bottomrule
\end{tabular}
\end{table}


\subsection{Probabilidad de denegacion en colas finitas}

En colas finitas, todos los valores de $\rho$ tienen referencia estacionaria porque el sistema bloquea arribos cuando alcanza su capacidad total. La Tabla~\ref{tab:mm1-denial-probability} compara la probabilidad de denegacion simulada con la teorica para cada capacidad de cola.

% Generado desde results/mm1/experiments/mm1_summary.csv
\begin{table}[htbp]
\centering
\caption{Probabilidad de denegacion simulada y teorica para colas finitas.}
\label{tab:mm1-denial-probability}
\begin{tabular}{lccccc}
\toprule
$\rho$ & $B=0$ & $B=2$ & $B=5$ & $B=10$ & $B=50$ \\
\midrule
0.25 & \shortstack{0.1994\\(0.2000)} & \shortstack{0.0118\\(0.0118)} & \shortstack{0.0002\\(0.0002)} & \shortstack{0.0000\\(0.0000)} & \shortstack{0.0000\\(0.0000)} \\
0.50 & \shortstack{0.3335\\(0.3333)} & \shortstack{0.0660\\(0.0667)} & \shortstack{0.0081\\(0.0079)} & \shortstack{0.0003\\(0.0002)} & \shortstack{0.0000\\(0.0000)} \\
0.75 & \shortstack{0.4288\\(0.4286)} & \shortstack{0.1549\\(0.1543)} & \shortstack{0.0519\\(0.0513)} & \shortstack{0.0113\\(0.0109)} & \shortstack{0.0000\\(0.0000)} \\
1.00 & \shortstack{0.5009\\(0.5000)} & \shortstack{0.2492\\(0.2500)} & \shortstack{0.1422\\(0.1429)} & \shortstack{0.0836\\(0.0833)} & \shortstack{0.0199\\(0.0192)} \\
1.25 & \shortstack{0.5556\\(0.5556)} & \shortstack{0.3379\\(0.3388)} & \shortstack{0.2533\\(0.2531)} & \shortstack{0.2146\\(0.2148)} & \shortstack{0.1999\\(0.2000)} \\
\bottomrule
\end{tabular}
\par\smallskip
\begin{minipage}{0.95\linewidth}
\centering\footnotesize
Cada celda muestra simulacion media y, entre parentesis, valor teorico.
\end{minipage}
\end{table}


\subsection{Graficos generados}

Los graficos originales se generaron en SVG y luego se copiaron al informe como PDF dentro de \texttt{report/figures/mm1/}. De esta manera se evita depender de soporte SVG adicional al compilar en Overleaf.

\begin{figure}[htbp]
    \centering
    \includegraphics[width=0.82\textwidth]{mm1_infinite_average_number_in_system.pdf}
    \caption{Promedio de clientes en el sistema para cola infinita estable. Solo se incluyen casos con $\rho < 1$, donde existe valor estacionario teorico.}
    \label{fig:mm1-infinite-number-system}
\end{figure}

\begin{figure}[htbp]
    \centering
    \includegraphics[width=0.82\textwidth]{mm1_infinite_average_number_in_queue.pdf}
    \caption{Promedio de clientes en cola para cola infinita estable. La comparacion con teoria se restringe a $\rho < 1$.}
    \label{fig:mm1-infinite-number-queue}
\end{figure}

\begin{figure}[htbp]
    \centering
    \includegraphics[width=0.82\textwidth]{mm1_infinite_average_time_in_system.pdf}
    \caption{Tiempo promedio en el sistema para cola infinita estable.}
    \label{fig:mm1-infinite-time-system}
\end{figure}

\begin{figure}[htbp]
    \centering
    \includegraphics[width=0.82\textwidth]{mm1_infinite_average_time_in_queue.pdf}
    \caption{Tiempo promedio de espera en cola para cola infinita estable.}
    \label{fig:mm1-infinite-time-queue}
\end{figure}

\begin{figure}[htbp]
    \centering
    \includegraphics[width=0.82\textwidth]{mm1_infinite_server_utilization.pdf}
    \caption{Utilizacion del servidor para cola infinita estable.}
    \label{fig:mm1-infinite-utilization}
\end{figure}

\begin{figure}[htbp]
    \centering
    \includegraphics[width=0.82\textwidth]{mm1_denial_probability_by_capacity.pdf}
    \caption{Probabilidad de denegacion de servicio en colas finitas para distintas capacidades de espera $B$.}
    \label{fig:mm1-denial-probability}
\end{figure}

\subsection{Evolucion temporal}

Ademas de las medidas finales, se genero una corrida representativa de cola infinita con $\rho=0.75$ y horizonte de $100$ horas para observar la evolucion del estado del sistema. Esta serie temporal no reemplaza las estimaciones finales de las tablas, que usan $10000$ horas y $10$ replicas, sino que muestra visualmente la dinamica de llegadas, esperas y salidas a lo largo del tiempo.

\begin{figure}[htbp]
    \centering
    \includegraphics[width=0.82\textwidth]{mm1_temporal_clientes_rho_075.pdf}
    \caption{Evolucion temporal de clientes en sistema y clientes en cola para una corrida representativa con $\rho=0.75$.}
    \label{fig:mm1-temporal-clientes}
\end{figure}

\section{Marco teorico del modelo de inventario}

El segundo modelo del trabajo practico representa un sistema de inventario con revision continua y politica $(Q,r)$. La variable principal es el inventario disponible $I_d$ al final del dia $d$. La politica queda definida por dos parametros: $Q$, cantidad fija a ordenar, y $r$, punto de reorden. Cuando la posicion de inventario cae a $r$ o menos, se emite una orden por $Q$ unidades.

La posicion de inventario se calcula como
\[
    PI_d = I_d + O_d,
\]
donde $O_d$ es la cantidad total pendiente de recepcion al momento de revisar la politica. Como en este modelo se adoptan ventas perdidas, la demanda no satisfecha no queda pendiente para dias posteriores. Por ese motivo no se resta una cartera de pedidos atrasados en la posicion de inventario.

La demanda diaria se modela como una variable aleatoria de Poisson:
\[
    D_d \sim \mathrm{Poisson}(\lambda_D),
\]
donde $\lambda_D$ representa la demanda media diaria. La demanda satisfecha del dia es $\min(I_d, D_d)$ y las unidades faltantes son la diferencia positiva entre demanda e inventario disponible.

El costo total de una corrida combina tres componentes:
\[
    C = K N + h \sum_{d=1}^{T} I_d + p \sum_{d=1}^{T} F_d,
\]
donde $K$ es el costo fijo de emitir una orden, $N$ es la cantidad de ordenes emitidas, $h$ es el costo de mantenimiento por unidad y por dia, $T$ es el horizonte de simulacion, $p$ es el costo por unidad faltante y $F_d$ son las unidades faltantes del dia $d$.

La medida de servicio usada en los resultados es
\[
    \text{nivel de servicio} =
    \frac{\text{demanda satisfecha}}{\text{demanda total}}.
\]
Esta metrica resume que proporcion de la demanda pudo atenderse directamente desde inventario disponible.

\subsection{Referencia esperada simplificada}

El modelo de inventario con demanda Poisson diaria, tiempo de entrega y ventas perdidas no tiene una expresion cerrada unica tan directa como el modelo M/M/1. Para contar con una referencia teorica esperada de comparacion, se usa una aproximacion clasica por ciclos de reposicion.

Sea $\lambda_D$ la demanda media diaria, $T$ el horizonte, $L_e$ el tiempo de entrega, $K_o$ el costo fijo de ordenar, $h$ el costo de mantenimiento y $p$ el costo de faltante. La demanda esperada del horizonte es
\[
    D_T = \lambda_D T,
\]
y la cantidad esperada de ciclos de reposicion se aproxima por
\[
    N_Q = \frac{D_T}{Q}.
\]

El inventario promedio esperado se aproxima mediante stock de ciclo mas stock de seguridad:
\[
    \bar I_{\mathrm{esp}} =
    \frac{Q}{2} + \max(r-\lambda_D L_e,0).
\]

Para medir el riesgo de faltantes durante el tiempo de entrega se define
\[
    D_{L_e} \sim \mathrm{Poisson}(\lambda_D L_e)
\]
y se estima el faltante esperado por ciclo como
\[
    E[(D_{L_e}-r)^+].
\]
Por lo tanto, el costo esperado simplificado usado como referencia es
\[
    C_{\mathrm{esp}} =
    K_o N_Q + hT\bar I_{\mathrm{esp}}
    + pN_Q E[(D_{L_e}-r)^+].
\]

Esta referencia no reemplaza a la simulacion: resume el sistema mediante demanda media, ciclos promedio y riesgo de faltante durante el tiempo de entrega. Sirve como punto de contraste para evaluar si la simulacion Python produce costos y decisiones de politica razonables.

\section{Metodologia del modelo de inventario}

La simulacion se implemento por periodos diarios, usando un horizonte de $365$ dias y $10$ replicas independientes por politica. La semilla base fue $20260626$ y las semillas de cada replica se derivaron de esa base para conservar reproducibilidad.

Los parametros comunes fueron:

\begin{itemize}
    \item Demanda diaria: Poisson con media de $20$ unidades por dia.
    \item Inventario inicial: $120$ unidades.
    \item Tiempo de entrega: $3$ dias.
    \item Costo fijo de orden: $120$.
    \item Costo de mantenimiento: $0.35$ por unidad por dia.
    \item Costo de faltante: $8$ por unidad faltante.
\end{itemize}

Se compararon tres politicas:

\begin{center}
\begin{tabular}{lrrl}
\toprule
Politica & $Q$ & $r$ & Interpretacion \\
\midrule
\texttt{bajo\_inventario} & 80 & 45 & Menor inventario promedio, mayor riesgo de faltantes \\
\texttt{balanceada} & 120 & 70 & Punto medio entre mantenimiento y faltantes \\
\texttt{alto\_servicio} & 160 & 95 & Mayor nivel de servicio, mayor mantenimiento \\
\bottomrule
\end{tabular}
\end{center}

En cada dia simulado se aplico la siguiente secuencia:

\begin{enumerate}
    \item Recibir ordenes pendientes cuyo dia de llegada sea el dia actual.
    \item Generar la demanda diaria.
    \item Satisfacer la demanda hasta agotar el inventario disponible.
    \item Registrar unidades faltantes cuando la demanda supera al inventario.
    \item Acumular costo de mantenimiento y costo de faltante.
    \item Calcular la posicion de inventario.
    \item Emitir una orden de cantidad $Q$ si la posicion es menor o igual que $r$.
    \item Acumular el costo fijo de orden y registrar el estado diario.
\end{enumerate}

El tratamiento de faltantes fue de ventas perdidas. Por lo tanto, una unidad no satisfecha se contabiliza como faltante y genera costo, pero no se arrastra como pedido pendiente al dia siguiente. Esta decision mantiene el modelo simple y consistente con la comparacion de politicas de servicio.

Las salidas de Python se guardaron en \texttt{results/inventory/experiments/}. El archivo \texttt{inventario\_corridas.csv} contiene una fila por replica, \texttt{inventario\_resumen.csv} resume cada politica y \texttt{inventario\_series.csv} conserva la evolucion diaria usada para los graficos.


\section{Resultados del modelo de inventario}

La Tabla~\ref{tab:inventario-resumen-politicas} resume el desempeno de las tres politicas. Con los parametros elegidos, la politica \texttt{balanceada} obtuvo el menor costo total promedio: $16304.39$, con un nivel de servicio de $99.96\%$ y un promedio de $2.70$ unidades faltantes por corrida. La politica \texttt{alto\_servicio} elimino los faltantes, pero su mayor inventario promedio incremento el costo de mantenimiento y llevo el costo total promedio a $20235.96$. La politica \texttt{bajo\_inventario} redujo el inventario disponible, pero acumulo $485.80$ unidades faltantes promedio y alcanzo un costo total promedio de $17998.25$.

% Generado desde results/inventory/experiments/inventario_resumen.csv
\begin{table}[htbp]
\centering
\caption{Resumen de desempeno por politica de inventario.}
\label{tab:inventario-resumen-politicas}
\begin{tabular}{lrrrrrr}
\toprule
Politica & $Q$ & $r$ & Costo total & IC 95\% & Servicio & Faltantes \\
\midrule
\texttt{bajo\_inventario} & 80 & 45 & 17998.25 & 234.94 & 93.37\% & 485.80 \\
\texttt{balanceada} & 120 & 70 & 16304.39 & 82.18 & 99.96\% & 2.70 \\
\texttt{alto\_servicio} & 160 & 95 & 20235.96 & 87.62 & 100.00\% & 0.00 \\
\bottomrule
\end{tabular}
\par\smallskip
\begin{minipage}{0.95\linewidth}
\centering\footnotesize
El nivel de servicio se calcula como demanda satisfecha sobre demanda total. Los faltantes corresponden a ventas perdidas.
\end{minipage}
\end{table}


La Tabla~\ref{tab:inventario-costos} separa los componentes de costo. En \texttt{bajo\_inventario}, el costo de faltante es relevante y compensa negativamente el ahorro de mantenimiento. En \texttt{alto\_servicio}, el costo de faltante desaparece, pero el costo de mantenimiento domina. La politica \texttt{balanceada} queda como recomendacion inicial porque reduce el costo total sin sacrificar de manera apreciable el nivel de servicio.

% Generado desde results/inventory/experiments/inventario_resumen.csv
\begin{table}[htbp]
\centering
\caption{Costos promedio por componente en el modelo de inventario.}
\label{tab:inventario-costos}
\begin{tabular}{lrrrr}
\toprule
Politica & Orden & Mantenimiento & Faltante & Total \\
\midrule
\texttt{bajo\_inventario} & 10188.00 & 3923.85 & 3886.40 & 17998.25 \\
\texttt{balanceada} & 7344.00 & 8938.79 & 21.60 & 16304.39 \\
\texttt{alto\_servicio} & 5508.00 & 14727.96 & 0.00 & 20235.96 \\
\bottomrule
\end{tabular}
\par\smallskip
\begin{minipage}{0.95\linewidth}
\centering\footnotesize
Todos los valores son promedios sobre las replicas de cada politica y usan el horizonte de 365 dias.
\end{minipage}
\end{table}


\subsection{Comparacion contra referencia esperada}

La Tabla~\ref{tab:inventario-comparacion-esperada} agrega una referencia esperada simplificada basada en demanda media, ciclos de reposicion y faltante esperado durante el tiempo de entrega. Esta referencia no busca reproducir exactamente cada detalle dinamico de la simulacion diaria, pero permite contrastar el orden de magnitud de los costos y el riesgo relativo de cada politica.

% Generado desde results/inventory/experiments/inventario_resumen.csv
\begin{table}[htbp]
\centering
\caption{Comparacion contra una referencia esperada simplificada de inventario.}
\label{tab:inventario-comparacion-esperada}
\begin{tabular}{lrrrrr}
\toprule
Politica & Costo esp. & Python & Dif. \% & Faltantes esp. & Faltantes Python \\
\midrule
\texttt{bajo\_inventario} & 27052.20 & 17998.25 & -33.47 & 1374.02 & 485.80 \\
\texttt{balanceada} & 16434.40 & 16304.39 & -0.79 & 23.99 & 2.70 \\
\texttt{alto\_servicio} & 20166.26 & 20235.96 & 0.35 & 0.00 & 0.00 \\
\bottomrule
\end{tabular}
\par\smallskip
\begin{minipage}{0.95\linewidth}
\centering\footnotesize
La referencia usa demanda media, ciclos esperados $D/Q$, inventario medio $Q/2+\max(r-\lambda_D L,0)$ y faltante esperado $E[(D_L-r)^+]$ con $D_L\sim Poisson(\lambda_D L)$.
\end{minipage}
\end{table}


La comparacion confirma el trade-off principal. La politica \texttt{bajo\_inventario} parece atractiva si solo se mira el costo de mantenimiento, pero su punto de reorden queda por debajo de la demanda media esperada durante el tiempo de entrega, por lo que aparece mayor riesgo de faltantes. En ese caso la referencia simplificada resulta conservadora, porque evalua el riesgo por ciclo y no representa con todo detalle la superposicion diaria de pedidos pendientes que si aparece en la simulacion. La politica \texttt{alto\_servicio} controla ese riesgo, aunque con mayor inventario promedio. La politica \texttt{balanceada} queda en una zona intermedia y por eso resulta la mejor alternativa en las corridas Python.

\subsection{Graficos generados}

Los graficos se generaron originalmente en SVG dentro de \texttt{figures/inventory/} y se copiaron como PDF dentro de \texttt{report/figures/inventory/} para facilitar la compilacion del informe en Overleaf.

\begin{figure}[htbp]
    \centering
    \includegraphics[width=0.82\textwidth]{inventario_disponible.pdf}
    \caption{Inventario disponible promedio por dia para cada politica.}
    \label{fig:inventario-disponible}
\end{figure}

\begin{figure}[htbp]
    \centering
    \includegraphics[width=0.82\textwidth]{demanda_diaria.pdf}
    \caption{Demanda diaria promedio observada en las replicas de cada politica.}
    \label{fig:inventario-demanda-diaria}
\end{figure}

\begin{figure}[htbp]
    \centering
    \includegraphics[width=0.82\textwidth]{costos_acumulados.pdf}
    \caption{Costo total acumulado promedio durante el horizonte simulado.}
    \label{fig:inventario-costos-acumulados}
\end{figure}

\begin{figure}[htbp]
    \centering
    \includegraphics[width=0.82\textwidth]{costo_total_por_politica.pdf}
    \caption{Costo total promedio por politica con intervalo de confianza del 95\%.}
    \label{fig:inventario-costo-total}
\end{figure}

\begin{figure}[htbp]
    \centering
    \includegraphics[width=0.82\textwidth]{nivel_servicio_por_politica.pdf}
    \caption{Nivel de servicio promedio por politica con intervalo de confianza del 95\%.}
    \label{fig:inventario-nivel-servicio}
\end{figure}


\clearpage
\section{Codigo y repositorio del proyecto}

El codigo fuente, los scripts de simulacion, los resultados exportados, las figuras y los archivos del informe se encuentran en el repositorio del proyecto:

\begin{center}
\href{https://github.com/gastonalt/Tp-3-simulaci-n}{github.com/gastonalt/Tp-3-simulaci-n}
\end{center}

Este enlace se registra como referencia provisoria del trabajo. Si el repositorio cambia en una etapa posterior, solo debe actualizarse esta seccion.


\end{document}
